% ===== §8 The keystone: the good that cannot be guaranteed =====
\section{The Unguaranteeability of the Good}
\label{p10:sec:keystone}

\noindent
The practice encountered a limit it could not cross from within: the practitioner, being a product of the cycle and possibly of an alienated one, cannot from the practitioner's own position certify the goodness of the practice. This section does not dissolve the limit. It establishes that the limit is the condition of the practice's being ethical, that the impossibility of self-certified goodness follows from the preceding development, and that this impossibility constitutes the good as a matter of ethics rather than calculation and as a process rather than a possessed state. This is the central juncture of the paper and of the series.

\subsection{The result stated}
\label{p10:sec:cruelty}

The result is stated in its direct form. No finite subject can guarantee that its own practice is good. The subject that loves cannot guarantee that its love does not constitute, or has not already constituted, an extraction; the subject that tends the field cannot certify, from within, that it tends rather than depletes; the subject that poses the two questions of \xref{p10:sec:criterion} in good faith cannot stand surety for the answers, the subject that poses them being the subject whose possible alienation is in question. There is no internal certificate of the subject's own goodness, no position within the subject from which the subject audits itself and pronounces its practice sound. The result is contrary to the demand, characteristic of the subject in a relation, for assurance that it does not injure the other, and the section establishes that the demand cannot be satisfied truthfully.

\subsection{The guarantee as a crossing of the is--ought gap}
\label{p10:sec:hume}

The result possesses a determinate form, the establishment of which is the work of the section, an analytic result supplying the rigour of a thesis otherwise statable only as a disposition. Consider the requirements on a subject that would guarantee the goodness of its own practice. To guarantee that an action will be good is to claim secure passage from a present fact, the state of the subject and the relation, to a normative outcome, the good the action will effect. It is to claim a valid inference from what is to what ought to follow. This is the inference Hume marked as admitting no valid passage: from premises concerning what is, no conclusion concerning what ought to be follows by any rule of inference; the passage from is to ought is a gap no finite reasoning traverses \parencite{hume2000}.

\begin{claim}[The guarantor of its own goodness as infinite reason]
To guarantee the goodness of one's own practice is to claim a secure inference from the present being of the subject and the relation to the good the action will realise. But the passage from is to ought is the gap Hume established that no finite reason traverses. A subject that could guarantee the goodness of its own practice would therefore be a subject that had traversed the is--ought gap, a subject for which the inference from fact to value is secure, which is an infinite reason: the term denoted, in the several traditions, as God, the Absolute, or the big Other this series has established does not exist. The demand for a guarantee of one's own goodness is, structurally, the demand to be that infinite reason. There being no such reason, no external standpoint, no metalanguage, no subject external to its own historical and relational production, there can be no guarantee. The impossibility of self-certified goodness is not a contingent limitation of human knowledge but the ethical consequence of finitude: of being a subject within the cycle rather than the infinite reason external to it.
\end{claim}

\noindent
This is the juncture at which the analytic and continental commitments of the paper coincide. The phenomenological result, that the subject cannot be assured that its love does not injure, acquires the rigour of a logical one: a subject claiming such assurance claims to have traversed Hume's gap, and thereby to be the infinite reason the series denies. The demand for a guarantee of one's own goodness is the covert invocation of the big Other, the term the series has established to be absent. The result is not a disposition but the form of finitude, with its demonstration.

\subsection{Unguaranteeability as the condition of ethics}
\label{p10:sec:reversal}

The consequence is the central thesis of the section. The unguaranteeability of the good is not the dissolution of ethics but its condition, that which constitutes the good a matter of ethics rather than of calculation.

\begin{claim}[Unguaranteeability as the condition of the ethical]
Were there a subject that could guarantee the goodness of its practice, that could compute, from present facts, the good its action would infallibly realise, ethics would be eliminated, and there would remain the execution of a determined optimal procedure: without risk, responsibility, trust, or exposure to the other. The ethical is constituted by the impossibility of the guarantee. Because the subject cannot guarantee that its love does not constitute extraction, the subject is committed to the sole course adequate to that impossibility: the continual sustaining of the question, the exposure of the practice to correction, and the consignment of the criterion of the subject's goodness, in part, to the other. The untraversable gap, far from licensing resignation, is the condition that directs the practice outward, and constitutes it as answerable and as ethical. A good admitting of guarantee would not be a good with respect to which the subject is required to be ethical; it would be a determined result.
\end{claim}

\noindent
The section thereby derives, from the result, the form the practice must assume, as its direct consequence rather than as a consolation. Since no subject certifies its own goodness from within, the good is not the possession or pronouncement of a single subject. It is polyphonic: consigned to the other, exposed to the other's correction, practised in a continual reciprocal disclosure in which each party is the other's sole available instrument of correction. This is the ground on which the preservation of divergence (\xref{p10:sec:subtractive}) is the practitioner's sole means of correcting the practitioner's own alienation: the subject not detecting its own displacement from within, the other party's capacity to register extraction is the sole instrument by which it is detected. The good is practised in the interval between two parties that each decline the guarantee and submit to correction.

\begin{claim}[The private certainty of one's own goodness as the cardinal error]
The single subject that attains, in isolation, certainty that the relation is good has, in that certainty, committed the cardinal error: the possession of a determination that, on pain of invoking the infinite reason, no single party may possess. The moment of private certainty of one's own goodness is the moment of greatest hazard.
\end{claim}

\subsection{The good as unfinished process}
\label{p10:sec:unfinished}

A final consequence returns to the dialectical materialism of \xref{p10:sec:dialmat}. That the good admits of no guarantee entails that the good is not a state, attained, certified, and subsequently possessed. It is a process: reproduced recurrently, and at each reproduction liable to failure. A good admitting of guarantee would be a good settled, closed, and completed, and a settled good is, by the development of the paper, a good cashed, possessed, and thereby alienated. It is in virtue of admitting no guarantee that the good remains a process rather than a fixed doctrine; the unguaranteeability of the good is the condition of its not being extinguished. A good rendered certain would, in the rendering, be extinguished; the good subsists only insofar as it is non-possessible, only as a process to be traversed anew and liable always to failure. The sustainability of the good established in the prior paper and its unguaranteeability are one result: the good subsists as a living spiral precisely in virtue of admitting no settlement, once and for all, into a fixed circle of certainty.

The result, fully developed, possesses two aspects which are one. The unguaranteeability of the good is not its defect; the defect would be the demand for the guarantee, which is the covert invocation of the infinite reason that does not exist, and is thereby the inception of alienation. To demand assurance of one's own goodness is to claim the external standpoint, the metalanguage, the big Other, to decline, finally, the finitude that is the condition of being a subject within the cycle. The practice of the good is the renunciation of that demand: the acceptance of an unguaranteeable, defeasible, recurrently unfinished good, consigned reciprocally between two parties that each decline to be the party that is certain. That renunciation, of certainty, of the guarantee, of the position external to the cycle, is the developed content of \emph{wu wei} for two parties, and the central juncture on which the development of the paper rests. It is, moreover, one of two renunciations the paper requires, the renunciation of the place of the Other in its two forms: the renunciation of the place of the Other that would guarantee one's goodness, developed here, and the renunciation of the place of the Other that would possess the relation's generativity, developed as the avoidance of symbolic exploitation in \xref{p10:sec:exploitation}. The two are one renunciation under two descriptions, and their unity is the content of the practice of the good. The juncture does not bear the weight of a guarantee, and is not required to. It bears the weight of a good that subsists precisely in virtue of admitting no certainty.
